\clearpage
\begingroup
\HistoricalHead{248}{L}{THEORETISCHE ASTRONOMIE. NACHLASS.}
\noindent\hspace*{\parindent}\textbf{Dritte Hypothese}\quad $\AstNode=80^0\,55'\,18''$\qquad $i=10^0\,36'\,30''$.
\par\medskip\noindent\hspace*{\parindent}[Die Rechnung ergab:]
{\small\setlength{\arraycolsep}{4pt}
\[
\begin{array}{rccc}
&\text{Erste Beobachtung}&\text{[Zweite Beobachtung]}&\text{[Dritte Beobachtung]}\\
{}[\log\Delta]&0{,}292\,9753&0{,}335\,8423&0{,}384\,6397\\
{}[v]&68^0\,38'\,45''31&73^0\,10'\,53''70&77^0\,58'\,31''18\\
{}[\log r]&0{,}435\,5885&0{,}431\,1483&0{,}427\,9709.
\end{array}
\]}
\noindent{\small [S. 7]}\par
\noindent\hspace*{\parindent}[Die] \textls[100]{Berechnung der Bahn aus den äussern Beobachtungen} [und die Berechnung der mittleren Beobachtung ergab
\[
\text{Fehler in }\log r'=+6639\qquad\text{Fehler in }v'=-99''45.]
\]
\begin{center}\textls[100]{Fehler der drei Hypothesen.}\end{center}
\noindent{\small [S. 8 u. 11]}
\[
\begin{array}{ccc|rr}
\multicolumn{3}{c}{}&[\Delta r]&[\Delta v]\\
\mathrm{III}&\AstNode&i&+6639&-99''45\\
\mathrm{I}&\AstNode+437''&i&-1146&-18''05\\
\mathrm{II}&\AstNode&i+176''&+2041&-15''375.
\end{array}
\]
\noindent\hspace*{\parindent}[Aus den Fehlern der Hypothesen III und I schliesst man:
\[
\frac{\partial\Delta r}{\partial\AstNode}=-\frac{7785}{437}\qquad\text{und}\qquad
\frac{\partial\Delta v}{\partial\AstNode}=\frac{11750}{437}
\]
\noindent und aus denen der Hypothesen III und II:
\[
\frac{\partial\Delta r}{\partial i}=-\frac{4598}{176}\qquad\text{und}\qquad
\frac{\partial\Delta v}{\partial i}=\frac{8407{,}5}{176}.
\]
\noindent\hspace*{\parindent}Bestimmt man sodann die an die Hypothese III anzubringenden Correctionen $\Delta i$ und $\Delta\AstNode$ aus der Gleichung
\[
\begin{aligned}
\Delta r&=-\frac{\partial\Delta r}{\partial\AstNode}\Delta\AstNode-\frac{\partial\Delta r}{\partial i}\Delta i,\\
\Delta v&=-\frac{\partial\Delta v}{\partial\AstNode}\Delta\AstNode-\frac{\partial\Delta v}{\partial i}\Delta i,
\end{aligned}
\]
\clearpage
\HistoricalHead{249}{R}{BEMERKUNGEN.}
\SourceLeaf{254}
\noindent so wird diese Gleichung\textsuperscript{*)}]
\begingroup\renewcommand{\thefootnote}{[*)}\footnotetext{Vgl. \textit{Theoria motus}, art. 120—122.]}\endgroup
\[
\begin{aligned}
6639&=7785x+4598y,\\
9945&=11750x+8407{,}5y,
\end{aligned}
\]
\noindent [wo $x=\dfrac{\Delta\AstNode}{437''}$ und $y=\dfrac{\Delta i}{176''}$ gesetzt ist.

Die Auflösung ergibt]
\[
\log x=9{,}946\,0893\qquad[\log y=8{,}7125n],
\]
\noindent [also
\[
\Delta\AstNode=]385''99\qquad[\Delta i=-9''08].
\]
\noindent [Es ergibt sich]
\[
\begin{array}{rr}
&80^0\,55'\,18''\\
&6\;26\\\cline{2-2}
{}[\AstNode=]&81\;1\;44
\end{array}
\qquad
\begin{array}{rr}
&10^0\,36'\,30''\\
&-9''08\\\cline{2-2}
{}[i=]&10\;36\;20''92.
\end{array}
\]
\noindent{\small [S. 12 u. 13]}\par
\noindent\hspace*{\parindent}[Mit diesen Werten als] \textbf{Vierte Hypothese} [wurde die] \textls[100]{Berechnung der Bahn} [von neuem ausgeführt und die III. Elemente gefunden\textsuperscript{**)}.]
\begingroup\renewcommand{\thefootnote}{[**)}\footnotetext{Vgl. Werke VI, S. 201.]}\endgroup
\clearpage
\endgroup
